\documentclass[a4paper]{report}
\usepackage{setspace}

\pagestyle{plain}
\usepackage{amssymb,graphicx,color}
\usepackage{amsfonts}
\usepackage{booktabs}
\usepackage[section]{placeins}
\usepackage{latexsym}
\usepackage{amsmath}
\usepackage{float}
\usepackage{multirow}
\usepackage{siunitx}
\usepackage[a4paper, margin = 3cm, bottom = 2.5cm]{geometry}

\newtheorem{theorem}{THEOREM}
\newtheorem{lemma}[theorem]{LEMMA}
\newtheorem{corollary}[theorem]{COROLLARY}
\newtheorem{proposition}[theorem]{PROPOSITION}
\newtheorem{remark}[theorem]{REMARK}
\newtheorem{definition}[theorem]{DEFINITION}
\newtheorem{fact}[theorem]{FACT}

\newtheorem{problem}[theorem]{PROBLEM}
\newtheorem{exercise}[theorem]{EXERCISE}
\def \set#1{\{#1\} }

\newenvironment{proof}{%
  PROOF: \begin{quotation}%
}{%
  $\Box$ \end{quotation}%
}

\newcommand{\nats}{\mbox{\( \mathbb N \)}}
\newcommand{\rat}{\mbox{\(\mathbb Q\)}}
\newcommand{\rats}{\mbox{\(\mathbb Q\)}}
\newcommand{\reals}{\mbox{\(\mathbb R\)}}
\newcommand{\ints}{\mbox{\(\mathbb Z\)}}

%%%%%%%%%%%%%%%%%%%%%%%%%%
\renewcommand{\thefootnote}{}


\title{{\vspace{-14em}}%
{{\Huge Cointegration-Based Spread Prediction for Mean-Variance Portfolio Optimisation }}\\
{\large Evaluating whether spread prediction between cointegrated asset pairs produces more reliable return estimates than historical means for Markowitz optimisation }\\
}
\date{Submission date: 24th April 2026}
\author{Declan Loo\thanks{
{\bf Disclaimer:}
This report is submitted as part requirement for the BSc Degree in Computer Science at UCL. It is
substantially the result of my own work except where explicitly indicated in the text. The report may be freely copied and distributed provided the source is explicitly acknowledged.
}%
\\ \\
BSc Computer Science\\ \\
Supervisor's name: Dr. Tomaso Aste}


\begin{document}

\onehalfspacing
\maketitle
\begin{abstract}
Summarise your report concisely.
\end{abstract}

\tableofcontents
\setcounter{page}{1}


\chapter{Introduction}
This is just a bare minimum to get started.  There is unlimited guidance on using latex, e.g., {\tt https://en.wikibooks.org/wiki/LaTeX}.   You are still responsible to check the detailed requirements of a project, including formatting instructions, see \\
{\tt https://moodle.ucl.ac.uk/pluginfile.php/3591429/mod\_resource/content/7/UGProjects2017.pdf}.
Leave at least a line of white space when you want to start a new paragraph.

\section{Section 1}
Chapters should contain numbered sections and sub-sections.

\subsection{Sub-Section 1}
Remember that this is a structured technical/academic report, not an essay.

\section{Mathematical Notation}
Mathematical expressions are placed inline between dollar signs, e.g. $\sqrt 2, \sum_{i=0}^nf(i)$, or in display mode
\[ e^{i\pi}=-1\] and another way, this time with labels,
\begin{align}
\label{line1} A=B\wedge B=C&\rightarrow A=C\\
&\rightarrow C=A\\
\intertext{note that}
n!&=\prod_{1\leq i\leq n}i \\
\int_{x=1}^y \frac 1 x \mathrm{d}x&=\log y
\end{align}
We can refer to labels like this \eqref{line1}.

\chapter{Background and Literature Review}

The main chapters should start with a brief summary -- this chapter is about, or introduces, some aspect of the work.

\section{Main Content}
Often lots of citations here (and elsewhere), e.g. \cite{Rey:D} or \cite[Theorem 2.3]{PriorNOP70}.   Bibtex can help with this, but is not essential. If you want pictures, try

\begin{center}
\end{center}
You can use
\begin{itemize}
\item lists
\item like this
\end{itemize}
or numbered
\begin{enumerate}
\item like this,
\item or this
\end{enumerate}
but don't overdo it.

\section{More Content}

\section{Summary}
The main chapters should have a short summary, highlighting what the chapter has covered and any key points that need emphasising.

\chapter{Analysis}
If you have a formal theorem you might try this.
\begin{definition}\label{def}
See definition~\ref{def}.
\end{definition}
\begin{theorem}
For all $n\in\nats,\; 1^n=1$.
\end{theorem}
\begin{proof}
By induction over $n$.
\end{proof}

\section{Methodology}

\section{Implementation Details}

\chapter{Testing and Results}

\section{Testing Methodology}

\section{Results}

\subsection{Cointegration Screening}

\begin{table}[htbp]
\centering
\caption{Engle-Granger Cointegration Screening Results (9 Candidate Pairs, In-Sample: Jan 2019 -- Dec 2023)}
\label{tab:cointegration_screening}
\resizebox{\textwidth}{!}{%
\begin{tabular}{l l c c c c c c c}
\toprule
Dependent & Independent & Hedge Ratio ($\hat{\beta}$) & Intercept ($\hat{\alpha}$) & ADF Stat & $p$-value & $y$ is $I(1)$ & $x$ is $I(1)$ & Cointegrated \\
\midrule
GS.N  & MS.N   & 0.780 &  2.351 & $-$3.667 & 0.020 & Yes & Yes & \textbf{Yes} \\
KO.N  & PEP.O  & 0.641 &  0.789 & $-$3.470 & 0.035 & Yes & Yes & \textbf{Yes} \\
CVX.N & XOM.N  & 0.698 &  1.813 & $-$3.234 & 0.064 & Yes & Yes & No \\
AMD.O & NVDA.O & 0.811 &  1.999 & $-$2.310 & 0.368 & Yes & Yes & No \\
AMZN.O & MSFT.O & 0.516 & 2.027 & $-$2.141 & 0.455 & Yes & Yes & No \\
JNJ.N & PFE.N  & 0.423 &  3.466 & $-$1.908 & 0.576 & Yes & Yes & No \\
META.O & GOOGL.O & 0.543 & 2.953 & $-$1.681 & 0.685 & Yes & Yes & No \\
TSM.N  & NVDA.O  & 0.476 & 3.071 & $-$1.340 & 0.817 & Yes & Yes & No \\
BAC.N  & JPM.N   & 0.854 & $-$0.666 & $-$1.230 & 0.850 & Yes & Yes & No \\
\bottomrule
\multicolumn{9}{l}{\footnotesize Augmented Engle-Granger test with MacKinnon (1994) critical values at 5\% significance ($p<0.05$). Log prices used throughout.} \\
\multicolumn{9}{l}{\footnotesize $I(1)$ verified by ADF on levels (non-stationary) and first differences (stationary). Sorted by $p$-value.} \\
\multicolumn{9}{l}{\footnotesize Source: Author's calculations on \texttt{01\_spread\_exploration.ipynb}.}
\end{tabular}%
}
\end{table}

\subsection{Spread Characterisation}

Based on the results in Table~\ref{tab:cointegration_screening}, two pairs were confirmed cointegrated at the 5\% significance level: GS.N/MS.N and KO.N/PEP.O. CVX.N/XOM.N was borderline ($p = 0.064$) and excluded. Spreads for each confirmed pair are computed as the OLS residual $S_t = \ln P^y_t - \hat{\beta}\ln P^x_t - \hat{\alpha}$ and characterised by their half-life, Hurst exponent, z-score, and spread statistics.

Table~\ref{tab:spread_characterisation} and Figures~\ref{fig:normalised_prices}--\ref{fig:spread_rolling_zscore}
present the spread characterisation results for the two cointegrated pairs.

\begin{table}[htbp]
\centering
\caption{Spread Characterisation Summary --- Cointegrated Pairs (In-Sample: Jan 2019 -- Dec 2023)}
\label{tab:spread_characterisation}
\resizebox{\textwidth}{!}{%
\begin{tabular}{l l c c c c c c c}
\toprule
Dependent & Independent & Half-Life (days) & Hurst Exponent & Current $z$-Score & Spread Mean ($\mu_S$) & Spread Std ($\sigma_S$) & ADF Stat & ADF $p$-value \\
\midrule
GS.N  & MS.N  & 41.8 & 0.403 & 0.675 & 0.000 & 0.056 & $-$3.666 & 0.005 \\
KO.N  & PEP.O & 40.3 & 0.421 & 0.809 & 0.000 & 0.047 & $-$3.469 & 0.009 \\
\bottomrule
\multicolumn{9}{l}{\footnotesize Half-life estimated via AR(1)/OU regression: $\kappa = \ln(2)/\text{half-life}$ gives $\kappa \approx 0.0166$ (GS/MS) and $0.0172$ (KO/PEP).} \\
\multicolumn{9}{l}{\footnotesize Hurst exponent $H < 0.5$ indicates mean-reverting (anti-persistent) dynamics. $z$-score uses 60-day rolling window at end of training period.} \\
\multicolumn{9}{l}{\footnotesize Spread mean is zero by construction (demeaned OLS residuals). Source: Author's calculations on \texttt{01\_spread\_exploration.ipynb}.}
\end{tabular}%
}
\end{table}

\begin{figure}[htbp]
    \centering
    \includegraphics[width=\textwidth]{figures/01_normalised_prices.pdf}
    \caption{Normalised log price series for the two cointegrated pairs over
    the in-sample period (January 2019 -- December 2023). Each series is
    normalised to zero at the start of the sample period to highlight co-movement.
    (a)~KO/PEP (Consumer Staples); (b)~GS/MS (Financials).}
    \label{fig:normalised_prices}
\end{figure}

\begin{figure}[htbp]
    \centering
    \includegraphics[width=0.9\textwidth]{figures/01_spread_characterisation.pdf}
    \caption{Spread characterisation summary for the two cointegrated pairs
    (January 2019 -- December 2023).
    (a) Half-life of mean reversion in trading days; dashed line marks the 60-day threshold.
    (b) Hurst exponent; dashed line marks the random walk boundary of 0.5.
    (c) Current $z$-score at the end of the training window; dotted lines mark $\pm2\sigma$.
    (d) ADF $p$-value on the spread; dashed line marks the 5\% significance level.
    Hatched bars indicate pairs that fail the respective threshold.}
    \label{fig:spread_characterisation}
\end{figure}

\begin{figure}[htbp]
    \centering
    \includegraphics[width=0.9\textwidth]{figures/01_spread_levels.pdf}
    \caption{Spread levels with $\pm1\sigma$ and $\pm2\sigma$ bands for each
    cointegrated pair over the in-sample period (January 2019 -- December 2023).
    The spread is defined as $S_t = \ln P^y_t - \hat{\beta}\ln P^x_t - \hat{\alpha}$.
    (a)~KO/PEP; (b)~GS/MS.}
    \label{fig:spread_levels}
\end{figure}

\begin{figure}[htbp]
    \centering
    \includegraphics[width=0.9\textwidth]{figures/01_spread_zscores.pdf}
    \caption{Rolling $z$-scores (60-day window) for the two cointegrated pairs
    over the in-sample period (January 2019 -- December 2023). Dotted lines mark
    the $\pm2\sigma$ entry thresholds. The $z$-score at the end of the training
    window feeds directly into the OU-implied return estimation.
    (a)~KO/PEP; (b)~GS/MS.}
    \label{fig:spread_zscores}
\end{figure}

\begin{figure}[htbp]
    \centering
    \includegraphics[width=0.9\textwidth]{figures/01_spread_distributions.pdf}
    \caption{Empirical spread distributions with overlaid normal density fits
    for the two cointegrated pairs (January 2019 -- December 2023).
    The spread mean is zero by construction as the OLS residual includes an intercept.
    (a)~KO/PEP; (b)~GS/MS.}
    \label{fig:spread_distributions}
\end{figure}

\begin{figure}[htbp]
    \centering
    \includegraphics[width=0.9\textwidth]{figures/01_spread_rolling_halflife.pdf}
    \caption{Rolling half-life of mean reversion (252-day window, EWM 21-day smoothed).
    Dashed line denotes the 60-day tradeability threshold.
    Shaded regions indicate periods where the smoothed half-life exceeded the threshold.
    (a)~KO/PEP; (b)~GS/MS.}
    \label{fig:spread_rolling_halflife}
\end{figure}

\begin{figure}[htbp]
    \centering
    \includegraphics[width=0.9\textwidth]{figures/01_spread_rolling_zscore.pdf}
    \caption{Rolling 60-day $z$-score for each cointegrated pair over the
    in-sample period (January 2019 -- December 2023). Dashed and dotted
    lines mark the $\pm1\sigma$ and $\pm2\sigma$ thresholds respectively;
    shaded regions indicate periods where the $z$-score exceeds $\pm2\sigma$,
    corresponding to potential entry signals.
    (a)~KO/PEP; (b)~GS/MS.}
    \label{fig:spread_rolling_zscore}
\end{figure}

\FloatBarrier

\subsection{Cointegration Validation}

Out-of-sample spreads are constructed using hedge ratios and intercepts
estimated solely on the in-sample period, ensuring no forward-looking
information is incorporated. Table~\ref{tab:oos_spread_stats} presents
the distributional shift in each spread across the IS and OOS periods,
while Table~\ref{tab:cointegration_validation} reports the formal
statistical test results. Figure~\ref{fig:oos_eg_pvalues} presents
the full IS vs OOS Engle-Granger $p$-value comparison across all
nine candidate pairs.

\begin{table}[htbp]
    \centering
    \caption{Spread statistics across in-sample (IS: January 2019 -- December
    2023) and out-of-sample (OOS: January 2024 -- December 2025) periods,
    constructed using in-sample hedge ratios.}
    \label{tab:oos_spread_stats}
    \resizebox{\textwidth}{!}{%
    \begin{tabular}{lcccccccc}
        \toprule
        & \multicolumn{2}{c}{Mean} & \multicolumn{2}{c}{Std Dev} & \multicolumn{2}{c}{Min} & \multicolumn{2}{c}{Max} \\
        \cmidrule(lr){2-3} \cmidrule(lr){4-5} \cmidrule(lr){6-7} \cmidrule(lr){8-9}
        Pair & IS & OOS & IS & OOS & IS & OOS & IS & OOS \\
        \midrule
        GS/MS   & 0.000 & $+$0.250 & 0.056 & 0.082 & $-$0.145 & $+$0.067 & $+$0.110 & $+$0.411 \\
        KO/PEP  & 0.000 & $+$0.170 & 0.047 & 0.110 & $-$0.173 & $-$0.016 & $+$0.111 & $+$0.367 \\
        \bottomrule
        \multicolumn{9}{l}{\footnotesize GS/MS OOS minimum of $+0.067$ lies entirely above the IS mean, indicating the spread never revisited its historical equilibrium.} \\
        \multicolumn{9}{l}{\footnotesize Source: Author's calculations on \texttt{02\_cointegration\_validation.ipynb}.}
    \end{tabular}%
    }
\end{table}

\begin{table}[htbp]
    \centering
    \caption{Cointegration validation statistics across IS (January 2019 --
    December 2023) and OOS (January 2024 -- December 2025) periods. ADF
    $p$-values below 0.05 indicate spread stationarity at the 5\% significance
    level. Half-life below 60 trading days and Hurst exponent below 0.5
    indicate mean reversion.}
    \label{tab:cointegration_validation}
    \resizebox{\textwidth}{!}{%
    \begin{tabular}{lccccccc}
        \toprule
        & \multicolumn{2}{c}{ADF $p$-value}
        & \multicolumn{2}{c}{Half-Life (days)}
        & \multicolumn{2}{c}{Hurst Exponent}
        & OOS \\
        \cmidrule(lr){2-3} \cmidrule(lr){4-5} \cmidrule(lr){6-7}
        Pair & IS & OOS & IS & OOS & IS & OOS & $z$-score \\
        \midrule
        GS/MS   & 0.005 & 0.452 & 41.8 &  81.1 & 0.403 & 0.362 & 1.015 \\
        KO/PEP  & 0.009 & 0.622 & 40.3 & 156.4 & 0.421 & 0.575 & 0.604 \\
        \bottomrule
        \multicolumn{8}{l}{\footnotesize Neither spread remains stationary OOS using IS hedge ratios. KO/PEP Hurst crosses above 0.5, indicating transition to near-random-walk behaviour.} \\
        \multicolumn{8}{l}{\footnotesize Source: Author's calculations on \texttt{02\_cointegration\_validation.ipynb}.}
    \end{tabular}%
    }
\end{table}

\begin{figure}[htbp]
    \centering
    \includegraphics[width=\textwidth]{figures/02_oos_eg_pvalues.pdf}
    \caption{IS vs OOS Engle-Granger $p$-values for all nine candidate pairs.
    Hatched bars indicate pairs that fail the 5\% significance threshold.
    GS/MS is the only pair persistent across both periods. KO/PEP degrades
    OOS; NVDA/TSM appears cointegrated OOS but is invalidated by NVDA.O
    failing the $I(1)$ precondition due to explosive price growth in 2024--2025.}
    \label{fig:oos_eg_pvalues}
\end{figure}

\begin{figure}[htbp]
    \centering
    \includegraphics[width=\textwidth]{figures/02_spread_is_oos_concat.pdf}
    \caption{Spread levels across the full sample period (January 2019 --
    December 2025). The dashed vertical line marks the IS/OOS boundary
    (January 2024). Horizontal lines and shading denote IS $\pm1\sigma$
    and $\pm2\sigma$ bands; persistent OOS displacement beyond these bounds
    indicates a structural break in the cointegrating relationship.
    (a)~KO/PEP; (b)~GS/MS.}
    \label{fig:spread_is_oos_concat}
\end{figure}

\begin{figure}[htbp]
    \centering
    \includegraphics[width=\textwidth]{figures/02_spread_oos_levels.pdf}
    \caption{Spread levels over the out-of-sample period (January 2024 --
    December 2025), constructed using in-sample hedge ratios and intercepts.
    Horizontal lines denote the IS mean and $\pm2\sigma$ bands estimated
    over the training window.
    (a)~KO/PEP; (b)~GS/MS.}
    \label{fig:spread_oos_levels}
\end{figure}

\begin{figure}[htbp]
    \centering
    \includegraphics[width=\textwidth]{figures/02_zscore_is_oos.pdf}
    \caption{Rolling 60-day $z$-score across the full sample period
    (January 2019 -- December 2025). The dashed vertical line marks the
    IS/OOS boundary. Shaded regions indicate $|z|>2\sigma$ exceedances.
    Persistent non-zero $z$-scores in the OOS period are consistent with
    a shift in the spread equilibrium level.
    (a)~KO/PEP; (b)~GS/MS.}
    \label{fig:zscore_is_oos}
\end{figure}

\begin{figure}[htbp]
    \centering
    \includegraphics[width=\textwidth]{figures/02_rolling_halflife_is_oos.pdf}
    \caption{Rolling half-life of mean reversion across the full sample
    period (January 2019 -- December 2025). Faint lines show raw estimates;
    bold lines show 21-day exponentially weighted smoothed values. The dashed
    horizontal line marks the 60-day tradability threshold. The dashed vertical
    line marks the IS/OOS boundary.
    (a)~KO/PEP; (b)~GS/MS.}
    \label{fig:rolling_halflife_is_oos}
\end{figure}

\FloatBarrier

\subsection{Return Estimation}

OU-implied expected returns are derived from the end-of-IS $z$-score for each
confirmed pair. A positive $z$-score (spread above mean) implies the spread is
expected to fall, producing a negative OU-implied return; a negative $z$-score
implies the spread is expected to rise, producing a positive return.
Table~\ref{tab:return_estimation} presents the sanity check confirming directional
consistency between the $z$-score and OU-implied return signal.

\begin{table}[htbp]
\centering
\caption{End-of-sample $z$-score validation and OU-implied expected returns (IS period: Jan 2019 -- Dec 2023)}
\label{tab:return_estimation}
\resizebox{\textwidth}{!}{%
\begin{tabular}{lrrrll}
\toprule
Pair & End-of-IS $z$-Score & OU-Implied $\mu$ (ann.) & Historical $\mu$ (ann.) & Signal & Consistent \\
\midrule
GS/MS  & $+$1.739 & $-$10.69\% & GS: $+$11.9\%, MS: $+$15.5\% & Spread well above mean & Yes \\
KO/PEP & $-$0.759 & $+$1.93\%  & KO: $+$6.4\%,\ PEP: $+$8.3\% & Spread below mean      & Yes \\
\bottomrule
\multicolumn{6}{l}{\footnotesize OU-implied $\mu$ computed via rolling 60-day OU model, annualised at 252 trading days per year.} \\
\multicolumn{6}{l}{\footnotesize Historical $\mu$ is the annualised sample mean of daily returns over the IS period.} \\
\multicolumn{6}{l}{\footnotesize Source: Author's calculations on \texttt{03\_return\_estimation.ipynb}.}
\end{tabular}%
}
\end{table}

\begin{figure}[htbp]
    \centering
    \includegraphics[width=\textwidth]{figures/03_cumulative_returns.pdf}
    \caption{Cumulative returns of each cointegrated spread versus its constituent
    assets over the in-sample period (January 2019 -- December 2023). All series
    are indexed to 1.0 at the start. The spread return reflects the daily change in
    the OLS residual $S_t = \ln P^y_t - \hat{\beta}\ln P^x_t - \hat{\alpha}$.
    (a)~KO/PEP; (b)~GS/MS.}
    \label{fig:cumulative_returns_is}
\end{figure}

\begin{figure}[htbp]
    \centering
    \includegraphics[width=\textwidth]{figures/03_mu_comparison.pdf}
    \caption{Expected return estimates (annualised, IS 2019--2023).
    (a)~OU-implied $\mu$ by pair: GS/MS has a large negative implied return
    (spread well above equilibrium at end of IS); KO/PEP has a small positive
    implied return (spread slightly below equilibrium).
    (b)~Historical mean $\mu$ by ticker: all four assets exhibit positive
    historical means, producing a very different allocation under mean-variance
    optimisation.}
    \label{fig:mu_comparison}
\end{figure}

\begin{figure}[htbp]
    \centering
    \includegraphics[width=0.75\textwidth]{figures/03_correlation_heatmap.pdf}
    \caption{Asset return correlation matrix for the four constituent tickers
    (GS.N, KO.N, MS.N, PEP.O) estimated over the in-sample period
    (January 2019 -- December 2023) using Ledoit-Wolf shrinkage covariance.
    GS and MS exhibit the highest within-pair correlation, consistent with
    their financial sector co-movement.}
    \label{fig:correlation_heatmap}
\end{figure}

\FloatBarrier

\subsection{Backtesting Results}

Table~\ref{tab:backtest_results} reports out-of-sample portfolio performance metrics
over the test period (January 2024 -- December 2025). The OU-Implied MPT portfolio
operates in spread space (GS/MS and KO/PEP as assets) under minimum-variance
weights; the Historical MPT portfolio operates in asset space (GS.N, KO.N, MS.N,
PEP.O) under minimum-variance weights with historical mean returns.

\begin{table}[htbp]
    \centering
    \caption{Out-of-sample portfolio performance metrics (test period: January 2024 -- December 2025). Vol.\ Reduction is relative to the equal-weight baseline; negative values indicate higher volatility than equal-weight.}
    \label{tab:backtest_results}
    \resizebox{\textwidth}{!}{%
    \begin{tabular}{lrrrrr}
    \toprule
     & Sharpe Ratio & Max Drawdown & Total Return & Ann.\ Volatility & Vol.\ Reduction \\
    \midrule
    OU-Implied MPT   & 1.359 & $-$8.7\%  & $+$34.2\% & 9.8\%  & $+$0.5\% \\
    Historical MPT   & 0.335 & $-$14.3\% & $+$12.1\% & 14.1\% & $-$43.4\% \\
    Equal-Weight     & 1.363 & $-$8.4\%  & $+$34.5\% & 9.8\%  & --- \\
    Buy \& Hold      & 1.403 & $-$30.0\% & $+$108.0\%& 27.5\% & --- \\
    S\&P 500 (SPY)   & 1.110 & $-$18.9\% & $+$44.3\% & 16.0\% & --- \\
    \bottomrule
    \multicolumn{6}{l}{\footnotesize Minimum-variance weights used for both MPT portfolios. Risk-free rate: 2\% p.a.} \\
    \multicolumn{6}{l}{\footnotesize Source: Author's calculations on \texttt{04\_backtest\_results.ipynb}.}
    \end{tabular}%
    }
\end{table}

\begin{table}[htbp]
    \centering
    \caption{Per-pair $z$-score mean-reversion backtest metrics (OOS: January 2024 -- December 2025). Entry at $\pm2.0\sigma$, exit at $\pm0.5\sigma$.}
    \label{tab:perpair_backtest}
    \resizebox{\textwidth}{!}{%
    \begin{tabular}{lrrrrrrrr}
    \toprule
    Pair & Sharpe & Max DD & Total Return & Ann.\ Vol & Trades & Days in Position & Win Rate \\
    \midrule
    GS.N/MS.N   & $-$0.285 & $-$20.7\% & $-$8.4\% & 17.2\% & 9 & 283 & 50.2\% \\
    KO.N/PEP.O  & $-$0.665 & $-$15.1\% & $-$6.3\% &  7.5\% & 8 & 312 & 47.8\% \\
    \bottomrule
    \multicolumn{8}{l}{\footnotesize Negative Sharpe ratios are consistent with OOS cointegration breakdown identified in \texttt{02\_cointegration\_validation.ipynb}.} \\
    \multicolumn{8}{l}{\footnotesize Source: Author's calculations on \texttt{04\_backtest\_results.ipynb}.}
    \end{tabular}%
    }
\end{table}

\begin{figure}[htbp]
    \centering
    \includegraphics[width=\textwidth]{figures/04_cumulative_returns.pdf}
    \caption{Cumulative growth of \$1 invested at the start of the OOS period
    (January 2024 -- December 2025) for the OU-Implied MPT, Historical MPT,
    equal-weight, and S\&P 500 benchmarks. The OU-Implied and equal-weight
    portfolios track closely throughout; the Historical MPT underperforms
    substantially.}
    \label{fig:oos_cumulative_returns}
\end{figure}

\begin{figure}[htbp]
    \centering
    \includegraphics[width=\textwidth]{figures/04_drawdown.pdf}
    \caption{Peak-to-trough drawdown over the OOS period (January 2024 --
    December 2025). Lower (closer to zero) is better. The OU-Implied MPT
    achieves a maximum drawdown of $-8.7\%$, substantially lower than the
    Historical MPT ($-14.3\%$) and S\&P 500 ($-18.9\%$).}
    \label{fig:oos_drawdown}
\end{figure}

\begin{figure}[htbp]
    \centering
    \includegraphics[width=\textwidth]{figures/04_rolling_sharpe.pdf}
    \caption{Rolling 60-day Sharpe ratio over the OOS period (January 2024 --
    December 2025). Time-varying risk-adjusted performance computed over a
    trailing 60-day window. Periods above zero indicate positive risk-adjusted
    returns in that window.}
    \label{fig:oos_rolling_sharpe}
\end{figure}

\begin{figure}[htbp]
    \centering
    \includegraphics[width=\textwidth]{figures/04_rolling_volatility.pdf}
    \caption{Rolling 60-day annualised volatility over the OOS period
    (January 2024 -- December 2025). The OU-Implied MPT maintains
    substantially lower and more stable volatility than the Historical
    MPT throughout the test period.}
    \label{fig:oos_rolling_volatility}
\end{figure}

\begin{figure}[htbp]
    \centering
    \includegraphics[width=\textwidth]{figures/04_weights_bar.pdf}
    \caption{Minimum-variance portfolio weights for both approaches.
    (a)~OU-Implied MPT: 46.6\% GS/MS spread, 53.4\% KO/PEP spread.
    (b)~Historical MPT: 11.7\% GS.N, 48.7\% KO.N, 0.0\% MS.N, 39.7\% PEP.O.
    The historical allocation concentrates heavily in KO and PEP due to their
    lower variance, while zeroing out MS.N.}
    \label{fig:portfolio_weights}
\end{figure}

\begin{figure}[htbp]
    \centering
    \includegraphics[width=\textwidth]{figures/04_pair_cumulative_returns.pdf}
    \caption{Cumulative returns from the $z$-score mean-reversion backtest engine
    for each cointegrated pair over the OOS period (January 2024 -- December 2025).
    Both pairs generate negative cumulative returns, consistent with persistent
    spread drift rather than mean reversion under IS-calibrated thresholds.
    (a)~GS.N/MS.N; (b)~KO.N/PEP.O.}
    \label{fig:pair_cumulative_returns}
\end{figure}

\begin{figure}[htbp]
    \centering
    \includegraphics[width=\textwidth]{figures/04_pair_zscores.pdf}
    \caption{Rolling $z$-scores with entry ($\pm2.0\sigma$) and exit ($\pm0.5\sigma$)
    thresholds for each cointegrated pair over the OOS period. Trade entry and exit
    points are marked. Persistent positive $z$-scores reflect the level shift in both
    spreads identified in the validation stage.
    (a)~GS.N/MS.N; (b)~KO.N/PEP.O.}
    \label{fig:pair_zscores}
\end{figure}

\begin{figure}[htbp]
    \centering
    \includegraphics[width=0.75\textwidth]{figures/04_efficient_frontier.pdf}
    \caption{Mean-variance efficient frontiers for both the OU-Implied (spread
    space) and Historical (asset space) portfolios. The two frontiers are not
    directly comparable as they operate over different asset universes. The
    minimum-variance portfolios used in the backtest are marked on each frontier.}
    \label{fig:efficient_frontier}
\end{figure}

\begin{figure}[htbp]
    \centering
    \includegraphics[width=\textwidth]{figures/04_return_distribution.pdf}
    \caption{Histogram of daily portfolio returns for the OU-Implied and Historical
    MPT portfolios over the OOS period (January 2024 -- December 2025). The
    OU-Implied distribution is narrower and more concentrated around positive values,
    reflecting its lower volatility and higher mean return.
    (a)~OU-Implied MPT; (b)~Historical MPT.}
    \label{fig:return_distribution}
\end{figure}

\FloatBarrier

\chapter{Dashboard}

\chapter{Discussion}

\section{Interpretation of Results}

Table~\ref{tab:spread_characterisation} shows that both cointegrated pairs exhibit
half-lives between 40 and 42 trading days and Hurst exponents below 0.5, confirming
statistically significant mean-reverting behaviour over the in-sample period. ADF tests
on the spread residuals reject the unit root hypothesis at the 1\% significance level
for both pairs, providing strong empirical support for the spread-based return
estimation approach in-sample.

However, Table~\ref{tab:cointegration_validation} reveals that neither spread
remains stationary OOS when IS hedge ratios are applied, with ADF $p$-values rising
to 0.452 and 0.622. This OOS breakdown directly affects the reliability of IS-calibrated
return signals, as evidenced by the negative Sharpe ratios in the per-pair backtest
(Table~\ref{tab:perpair_backtest}). Despite this, the portfolio-level OU-Implied MPT
substantially outperforms the Historical MPT on all three evaluation metrics
(Table~\ref{tab:backtest_results}), suggesting that the structural properties of
the spread universe — low volatility and low market beta — provide resilience
beyond the precision of the OU return estimates alone.

\chapter{Conclusions and Evaluation}
\section{Achievements}
Summarise the achievements to confirm the project goals have been met.
\section{Evaluation}
Evaluation of the work (this may be in a separate chapter if there is substantial evaluation).

\section{Future Work}
How the project might be continued, but don't give the impression you ran out of time!

The project could have continued by adding a rebalance option which allows the functionality of producing a portfolio with percentage of each asset and capital allocated in a way to minimise the number of transactions needed, to minimise transactional cost.

\appendix
\addcontentsline{toc}{chapter}{Bibliography}
\begin{thebibliography}{HHM99}


\bibitem[Pri70]{PriorNOP70}  %%only an example
A.~Prior.
\newblock The notion of the present.
\newblock {\em Studium Generale}, 23:  245--248, 1970.


\bibitem[Rey97]{Rey:D}
M.~Reynolds.
\newblock A decidable temporal logic of parallelism.
\newblock {\em Notre Dame Journal of Formal Logic}, 38(3):  419--436,
  1997.
\end{thebibliography}
\newpage
\addcontentsline{toc}{chapter}{Appendices}

\chapter{Selected US Equities: Candidate Universe for Pair Selection}

\begin{table}[h!]
\centering
\caption{US Equity Candidate Universe ranked by Market Capitalisation}
\label{tab:universe_mcap}
\begin{tabular}{cllr}
\toprule
\textbf{Rank} & \textbf{Company} & \textbf{Symbol} & \textbf{Market Cap in USD (29 Dec 2023)} \\
\midrule
1 & Microsoft & MSFT.O & \$2,795B \\
2 & Alphabet & GOOGL.O & \$1,755B \\
3 & Amazon & AMZN.O & \$1,570B \\
4 & Nvidia & NVDA.O & \$1,223B \\
5 & Meta & META.O & \$910B \\
6 & TSMC & TSM.N & \$501B \\
7 & JPMorgan & JPM.N & \$492B \\
8 & ExxonMobil & XOM.N & \$400B \\
9 & J\&J & JNJ.N & \$377B \\
10 & Costco & COST.O & \$293B \\
11 & Chevron & CVX.N & \$281B \\
12 & BofA & BAC.N & \$266B \\
13 & Coca-Cola & KO.N & \$255B \\
14 & AMD & AMD.O & \$238B \\
15 & PepsiCo & PEP.O & \$234B \\
16 & Intel & INTC.O & \$212B \\
17 & Pfizer & PFE.N & \$163B \\
18 & Morgan Stanley & MS.N & \$153B \\
19 & Goldman Sachs & GS.N & \$126B \\
20 & Schlumberger & SLB.N & \$74B \\
21 & Target & TGT.N & \$66B \\
22 & Halliburton & HAL.N & \$32B \\
23 & Delta Airlines & DAL.N & \$26B \\
24 & United Airlines & UAL.O & \$14B \\
\bottomrule
\end{tabular}
\end{table}


\chapter{Design Diagram of the Dashboard with Excalidraw}

\chapter{Project Plan}

\chapter{Interim Report}


\end{document}
